\chapter{Remotes and pull requests}
\label{ch:remotes}

\chaptergoals{3 of 6 in Part~I.}
{\cref{ch:git,ch:branches}. A GitHub account for the second half of the
chapter, from \cref{sec:forks} on.}
{exchange commits with other copies of a repository (clone, fetch, pull,
push), work with a fork, and take a change through a ShipSoft pull request.}

\section{Remotes}

A \emph{remote} is another copy of the repository that you exchange commits
with. On GitHub it is usually called \texttt{origin}. Before going to GitHub,
it helps to see that nothing magic happens there: a remote can just as well
be a directory on your own disk.

\begin{exercise}{A server on your own machine}
\label{ex:localserver}
Play two people, Alice and Bob, sharing one ``server'' repository. Next to
\file{git-playground} (not inside it):
\begin{shell}
git init --bare server.git              # a repository without a working tree
git clone server.git alice
git clone server.git bob                # (both warn: the repository is empty)
cd alice
vim notes.md                            # one line: # Notes
git add notes.md
git commit -m "docs: start the notes"
git push -u origin main                 # Alice publishes main
cd ../bob
git pull origin main                    # Bob gets it
\end{shell}
Now Bob starts work on a branch while Alice publishes more:
\begin{shell}
git switch -c bob-idea                  # in bob/
vim idea.md                             # one line: an idea
git add idea.md && git commit -m "docs: add an idea"
cd ../alice
vim notes.md                            # add a second line: more
git commit -am "docs: more notes" && git push
cd ../bob
git fetch                               # download, change nothing
git log --oneline --graph --all         # origin/main is ahead of Bob's base
git rebase origin/main                  # replay Bob's idea on top of it
git push -u origin bob-idea
\end{shell}
\end{exercise}

These are all the remote operations you need:

\begin{description}[style=nextline]
\item[\cmd{git clone <address>}] copies a repository and remembers where it
  came from as the remote \texttt{origin}.
\item[\cmd{git fetch}] downloads new commits and updates the
  \emph{remote-tracking} branches such as \texttt{origin/main}. It never
  touches your own branches or files, so it is always safe; afterwards you
  can look before you act (\cmd{git log main..origin/main}).
\item[\cmd{git pull}] is \cmd{git fetch} followed by integrating the remote
  branch into your current branch: by rebasing, with the \cmd{pull.rebase}
  setting from \cref{ex:gitfirst}.
\item[\cmd{git push -u origin <branch>}] uploads your branch. \cmd{-u}
  remembers the pairing, so later a plain \cmd{git push} suffices.
\item[\cmd{git push -{}-force-with-lease}] uploads a branch you have
  rebased, replacing your old commits on the remote. It refuses if somebody
  else pushed to that branch in the meantime; never use plain
  \cmd{-{}-force}.
\end{description}

\section{Forks}
\label{sec:forks}

You cannot push to the ShipSoft repositories themselves. Like most
open-source projects, SHiP uses \emph{forks}: GitHub makes you a personal copy
of the repository, you push your branches there, and you ask for them to be
taken into the original with a pull request. So you work with two remotes.

\begin{figure}[ht]
\centering
\begin{tikzpicture}[
  box/.style={draw=shipnavy, fill=shipcream, rounded corners=2pt, minimum width=36mm,
    minimum height=11mm, align=center, font=\sffamily\small},
  move/.style={-{Stealth[length=2mm]}, line width=0.7pt, color=shipnavy},
  lab/.style={font=\ttfamily\scriptsize, fill=white, inner sep=1pt, align=center}]
  \node[box] (up) {\textbf{upstream}\\\scriptsize github.com/ShipSoft/aegir};
  \node[box, right=34mm of up] (fork) {\textbf{origin} (your fork)\\\scriptsize github.com/you/aegir};
  \node[box, below=16mm of $(up)!0.5!(fork)$] (local) {\textbf{your clone}\\\scriptsize on your machine};
  \draw[move, dashed] ([yshift=1.5mm]fork.west) -- node[lab, above=1pt] {pull request} ([yshift=1.5mm]up.east);
  \draw[move] ([yshift=-1.5mm]up.east) -- node[lab, below=1pt] {``Fork'' button} ([yshift=-1.5mm]fork.west);
  \draw[move] (up.south) -- node[lab, left=2pt] {git fetch upstream\\git pull upstream main} (local.north west);
  \draw[move] (local.north east) -- node[lab, right=2pt] {git push origin\\\textrm{\scriptsize (your branch)}} (fork.south);
\end{tikzpicture}
\caption{The fork workflow. You read from \texttt{upstream}, write to your
fork (\texttt{origin}), and ask for your branch to be taken into
\texttt{upstream} with a pull request.}
\label{fig:forks}
\end{figure}

\begin{shell}
git clone git@github.com:<you>/aegir.git        # your fork becomes "origin"
cd aegir
git remote add upstream git@github.com:ShipSoft/aegir.git
git remote -v                                    # lists both
\end{shell}

In \cref{ex:localserver} there was one remote; with a fork, you fetch and pull
from \texttt{upstream} and push to \texttt{origin}.

GitHub accepts SSH (\texttt{git@github.com:\ldots}) or HTTPS
(\texttt{https://github.com/\ldots}) addresses. This booklet uses SSH, which
needs a key on lxplus that your GitHub account knows.

\begin{exercise}{An SSH key for GitHub, once}
\label{ex:sshkey}
On lxplus:
\begin{shell}
ssh-keygen -t ed25519 -C "your.name@cern.ch"   # accept the file name; choose a passphrase
cat ~/.ssh/id_ed25519.pub                      # the public key: one line
\end{shell}
Copy that line into GitHub: \emph{Settings}, \emph{SSH and GPG keys},
\emph{New SSH key}. Then check:
\begin{shell}
ssh -T git@github.com          # "Hi <you>! You've successfully authenticated"
\end{shell}
Only the \texttt{.pub} file ever leaves lxplus; the other one is your secret.
\end{exercise}

\begin{organisernote}
Check once that outgoing SSH from lxplus to \texttt{github.com} works on the
tutorial day. If it does not, use HTTPS addresses and \cmd{gh auth login}
(the GitHub CLI, from a pixi environment) instead.
\end{organisernote}

\section{Pull requests the ShipSoft way}
\label{sec:pullrequests}

A pull request (PR) asks the maintainers of \texttt{upstream} to take the
commits of one of your branches. The rules below are the ones configured for
the ShipSoft repositories (in \repo{.github}, \file{repo-config/}) and
described in each package's \file{CONTRIBUTING.md}.

\begin{enumerate}
\item \textbf{Start from an up-to-date \texttt{main}}, on a new branch:
\begin{shell}
git switch main && git pull upstream main
git switch -c fix-field-translation
\end{shell}
\item \textbf{Commit} in small, self-contained steps with Conventional-Commit
  messages.
\item \textbf{Push} the branch to your fork and \textbf{open the PR} against
  \texttt{main} of the ShipSoft repository, on the page GitHub shows you after
  the push, or with \cmd{gh pr create}. Say what the change does, why, and how
  you tested it.
\item \textbf{Checks.} Continuous integration builds and tests the branch
  automatically. The required checks are called ``All checks passed'' and, in
  most repositories, ``Lint passed''. A red check is yours to fix: push more
  commits to the same branch and the PR updates. \Cref{ch:ci} shows how this
  works.
\item \textbf{Review.} At least one approving review is required, and every
  review conversation must be resolved. A new push dismisses earlier
  approvals, so the reviewer sees what is actually merged.
\item \textbf{Up to date.} If \texttt{main} moved on in the meantime, the PR
  must be brought up to date, by rebasing, exactly as Bob did in
  \cref{ex:localserver}:
\begin{shell}
git fetch upstream
git rebase upstream/main
# run the tests again
git push --force-with-lease
\end{shell}
\item \textbf{Merge.} An approved, green PR goes into the merge queue, which
  rebases it onto \texttt{main}, runs the checks once more, and only then
  lands it. Merge commits are not possible, and nobody, maintainers included,
  can force-push to or delete \texttt{main}.
\item \textbf{Afterwards} update your own copies and delete the branch:
\begin{shell}
git switch main && git pull upstream main && git push origin main
git branch -d fix-field-translation
\end{shell}
\end{enumerate}

\begin{organisernote}
Rules as configured in \repo{.github} on 15 October 2026
(\file{repo-config/ruleset-protect.json}, \file{ruleset-review.json},
\file{repo-settings.json}): no merge commits; squash and rebase allowed;
linear history; merge queue with method REBASE; strict required status
checks ``All checks passed'' (and ``Lint passed'' for tier ``full''); one
approval; stale approvals dismissed; review threads must be resolved;
deletion and non-fast-forward pushes to the default branch blocked.
Trout, \texttt{.github} and Documentation have no required checks; the
forks EventCalc and ROOTEGPythia6 are not managed. Check again before the
session.
\end{organisernote}

You will take a real pull request through all of these steps in
\cref{ch:team}. Two things in step~4 still need explaining first: the
environment the checks run in (\cref{ch:pixi}), and the checks themselves
(\cref{ch:ci}).

\section{Git cheat sheet}
\label{sec:gitcheatsheet}

\begin{center}\small
\begin{tabularx}{\linewidth}{>{\ttfamily}l L}
\toprule
\normalfont\bfseries Command & \bfseries What it does \\
\midrule
git status                    & what changed, what is staged, which branch \\
git diff [-{}-staged]         & changes not yet staged [already staged] \\
git add <file>                & stage a change for the next commit \\
git commit -m "type: summary" & record the staged changes \\
git log -{}-oneline -{}-graph -{}-all & the history as a graph \\
git restore <file>            & discard unstaged changes to a file \\
git switch -c <branch>        & create a branch and switch to it \\
git switch <branch>           & switch to an existing branch \\
git fetch upstream            & download, change nothing locally \\
git pull upstream main        & fetch and rebase the current branch onto it \\
git rebase upstream/main      & replay your branch on the latest main \\
git rebase -{}-continue / -{}-abort & after resolving a conflict / give up \\
git push -u origin <branch>   & publish a branch to your fork \\
git push -{}-force-with-lease & publish a rebased branch (yours only!) \\
\bottomrule
\end{tabularx}
\end{center}

\begin{levelcheck}
\begin{itemize}
  \can clone, fetch, pull and push, and say what each one changes;
  \can bring an out-of-date branch up to date with a rebase and push it
       safely;
  \can set up a fork with \texttt{origin} and \texttt{upstream};
  \can list the steps a ShipSoft pull request goes through before it is
       merged.
\end{itemize}
\end{levelcheck}

\begin{solution}
Every command sequence in \cref{ch:git,ch:branches,ch:remotes} and in
\cref{ch:team} is exercised by \file{tests/test\_git\_workflow.sh} against
local bare repositories standing in for GitHub, including Alice and Bob,
an out-of-date PR, a rebase onto \texttt{upstream/main}, a conflict resolved
with \cmd{git rebase -{}-continue}, and a merge queue that leaves no merge
commits.
\end{solution}
